\documentclass[11pt]{article}
\usepackage[a4paper,margin=22mm]{geometry}
\usepackage[T1]{fontenc}
\usepackage{lmodern}
\usepackage{amsmath}
\usepackage{microtype}
\usepackage{graphicx}
\usepackage{xcolor}
\usepackage{float}
\usepackage{fancyhdr}
\usepackage[hidelinks]{hyperref}

\definecolor{Accent}{HTML}{21918C}
\definecolor{Muted}{HTML}{5B6570}

\newcommand{\Faq}{\textit{F.~aquilonia}}
\newcommand{\Fpo}{\textit{F.~polyctena}}

%% One block per manuscript figure: image, caption, short Methods, optional Note.
\newcommand{\figblock}[5]{%
  % #1 = pdf file (in ../Figures/), #2 = title, #3 = caption, #4 = methods, #5 = note (may be empty)
  \section*{\color{Accent}#2}
  \begin{figure}[H]\centering
    \includegraphics[width=\linewidth]{../Figures/#1}
  \end{figure}
  {\small\textbf{Figure caption.} #3\par}
  \medskip
  {\small\textbf{\color{Muted}Materials \& Methods.} #4\par}
  \if\relax\detokenize{#5}\relax\else
    \medskip{\small\textbf{\color{Accent}Note.} #5\par}\fi
  \clearpage}

\pagestyle{fancy}\fancyhf{}
\lhead{\small\color{Muted}Manuscript figures --- captions \& methods}
\rhead{\small\color{Muted}Poikela et al.}
\cfoot{\small\thepage}
\renewcommand{\headrulewidth}{0.4pt}
\setlength{\parskip}{0.4em}\setlength{\parindent}{0pt}

\title{\vspace{-3em}\textbf{\color{Accent}Manuscript figures}\\[0.2em]
  \large Captions and short Materials \& Methods}
\author{\normalsize Poikela et al.}
\date{\normalsize\today}

\begin{document}
\maketitle
\vspace{-1.2em}
\noindent\small\textit{One block per manuscript figure: the figure, its caption, a short
Methods paragraph stating exactly what data and analysis go into it, and --- where useful ---
a short Note discussing a subtlety a reader may raise. Most figures are
generated by \texttt{module\_di25/Figures.R} (each \texttt{fig\_*()} section); the
ancestry-mosaic circos (Fig.~1) is a heavy standalone raster generated by
\texttt{module\_di25/R/diem\_circos\_compare.R}.
Species colours in the line figures: \Faq{} teal, \Fpo{} yellow.}
\clearpage

%% ==========================================================================
\figblock{diem_circos_compare_snp_vs_ldreduced.png}
{Figure 1 --- Ancestry mosaic before and after LD complexity reduction}
{% caption
\textbf{Genome-wide DIEM ancestry mosaic of the diagnostic markers, per individual,
before and after LD complexity reduction.} Circular genotype heatmaps of 195 individuals ---
165 hybrid workers and 30 parental reference individuals: each concentric ring is one
individual, and the rings are ordered once --- innermost = most \Faq{}, outermost = most
\Fpo{}, by the per-SNP mean hybrid index --- with that same order applied to both panels, so a
ring at a given radius is the same individual in \textbf{(a)} and \textbf{(b)}. Markers are
arranged by genomic position around the circle (chromosomes~1--27). Cell colour is DIEM
ancestry: purple = homozygous \Faq{}, teal = heterozygous, yellow = homozygous \Fpo{}; missing
genotypes are blank. \textbf{(a)} All 51{,}612 diagnostic SNPs. \textbf{(b)} The same
individuals and markers after LD complexity reduction to 11{,}052 units at a 5\,cM merge cap
(4{,}010 eMLG consensus units from clusters of $>2$ markers, plus 7{,}042 representative SNPs
from 1--2-marker clusters). The large-scale ancestry structure --- the per-individual
\Faq{}--\Fpo{} gradient and the chromosomal ancestry blocks --- is retained after the
$\sim$4.7-fold reduction, with the redundant within-block SNP replication collapsed.}
{% methods
\emph{Data.} The same DI\,$>-25$ diagnostic markers and LD-reduced units as Figure~2
(51{,}612 SNPs; 11{,}052 units), displayed across all 195 genotyped individuals --- 165
\Faq{}\,$\times$\,\Fpo{} hybrid workers from 20 populations and 30 parental references (15 per
species). Genotypes are shown in three states --- homozygous \Faq{}, heterozygous, homozygous
\Fpo{} --- with missing calls left blank.

\emph{Units (panel b).} Each LD-reduced unit was summarised across its member markers:
clusters of more than two markers by an expected multi-locus genotype (eMLG) consensus
dosage, 1--2-marker clusters by their representative SNP. Each unit was oriented to the
\Faq{} allele by the sign of the parental allele-frequency difference and then discretised to
the three-state ancestry code; units are placed at their representative marker's genomic
position.

\emph{Individual order.} Individuals were ordered once by their per-SNP mean hybrid index
(fraction \Faq{}) and that single order was applied to both panels, so LD reduction is the
only difference between a and b.

\emph{Rendering.} Both panels are drawn as polar rasters
(\texttt{render\_diem\_circos()}, \texttt{module\_di25/R/diem\_circos\_core.R}) and written as
one 5200\,$\times$\,2750-pixel, 300-dpi PNG by
\texttt{module\_di25/R/diem\_circos\_compare.R}. Being dense per-cell rasters
(individuals\,$\times$\,tens of thousands of loci), the panels are emitted as a
high-resolution bitmap rather than vector graphics.}
{% note
\emph{Colour scheme differs from the line figures.} In this ancestry heatmap teal marks
\emph{heterozygotes} and the \Faq{}/\Fpo{} homozygotes are purple/yellow (the three-state
DIEM coding). In Figure~2 the teal/yellow pair instead marks the \emph{direction} of
sorting (\Faq{}/\Fpo{}). The consistent anchor across figures is yellow\,=\,\Fpo{}.}

%% ==========================================================================
\figblock{fig_recomb_directional.pdf}
{Figure 2 --- Direction of ancestry sorting vs recombination}
{% caption
\textbf{Ancestry sorting toward each parental species as a function of recombination, and
its dependence on the sorting threshold.}
\textbf{(a)} Fraction of units sorted toward \Faq{} (teal) or
\Fpo{} (yellow) across recombination deciles (\emph{x}: median cM/Mb per decile), at two
sorting thresholds --- a unit is ``sorted'' when at least 60\,\% (left) or 80\,\% (right)
of populations are near-fixed. Solid lines, LD-reduced units (eMLG); dashed lines,
individual SNPs. Grey bars give the number of independent eMLG units per decile (right
axis). The three lowest-recombination deciles (shaded) contain only 8, 9 and 38 independent
units, so their per-decile \emph{fractions} are individually imprecise --- a couple of units
shift them substantially; they are shaded for that reason, not because the underlying signal is
doubtful. The directional contrast at low recombination is therefore quantified in panel b,
which enters each unit as one observation along the continuous gradient rather than reading a
fraction off these sparse bins. Error bars are Wilson 95\,\% confidence intervals, shown for the
eMLG estimates only (SNPs are not independent). \textbf{(b)} Recombination coefficient of
each direction's sorting at each threshold: the slope of a logistic regression of
$P(\text{sorted toward the species})$ on $\log_{10}$ recombination, fitted across all
LD-reduced units; points and 95\,\% intervals are from a genomic block bootstrap that
resamples whole chromosomes. Negative values indicate sorting concentrated at low
recombination. Both directions concentrate at low recombination, but \Fpo{}-directed sorting
does so more strongly: the paired block-bootstrap difference
in coefficient (\Faq{}\,$-$\,\Fpo{}) is positive at all four thresholds and its 95\,\% CI
excludes zero at three of them ($\tau=0.6$: $+1.1$, 95\,\% CI $[0.05,\,2.2]$, two-sided
$p=0.04$; significant at $\tau=0.5$--$0.7$, marginal at $\tau=0.8$ where $p=0.06$, shown above
each threshold in panel b). Fitting the gradient this way is what makes
the low-recombination contrast estimable: the sparse bins of panel a cannot support a stable
per-decile fraction, whereas the regression enters every unit --- including the few large
\Fpo{}-sorted low-recombination blocks --- as one independent observation, with the
chromosome-level block bootstrap accounting for their non-independence.}
{% methods
\emph{Scope.} The entire pipeline underlying this figure --- LD-decay estimation, LD-based
complexity reduction, and the ancestry-sorting analysis --- is based exclusively on the
high-diagnostic-index (DI $>-25$) markers; other analyses in this study use the full marker
set.

%% TODO: 164 -> 165 hybrids once the DI25 clustering + sorting are rerun on the corrected
%% 195-individual sample (one hybrid was missing from the earlier sample info). Fig. 1
%% (circos) already displays the full 195; this analysis still reflects the 164-hybrid run.
\emph{Diagnostic markers and populations.} We analysed the 51{,}612 species-diagnostic
SNPs (diagnostic index $>-25$) genotyped in 164 \Faq{}\,$\times$\,\Fpo{} hybrid workers
from 20 populations, with 30 parental reference individuals, coded as 0/1/2 allele dosage.

\emph{LD-reduced units.} Linkage-disequilibrium clustering was rebuilt from scratch on the
diagnostic markers only, using hybrids alone (including parents would inflate LD through
parent--hybrid structure): chromosome-specific LD-decay, a per-marker local-LD statistic,
and a two-stage, quality-gated complexity reduction with a 0.5\,cM Stage-1 threshold and a
5\,cM merge cap. This collapsed the 51{,}612 SNPs to 11{,}052 LD-reduced units; clusters of
more than two markers are represented by an expected multi-locus genotype (eMLG) consensus,
and 1--2-marker clusters by their representative SNP.

\emph{Ancestry sorting.} Each unit was oriented to the \Faq{} allele by the sign of the
parental allele-frequency difference. A population was scored near-fixed toward a species
when its oriented allele frequency reached $\varphi=0.85$ (\Faq{}) or fell to $\le 1-\varphi$
(\Fpo{}). A unit was called sorted when the proportion of populations near-fixed reached the
threshold $\tau$; among sorted units the direction was assigned to the majority species when
a two-sided binomial test of $n_{\text{aqu}}$ among $n_{\text{fixed}}$ ($\alpha=0.05$)
rejected parity, and left unresolved otherwise. The only diagnosticity gate is the
DI\,$>-25$ marker selection itself; unlike the genome-wide Module-A sorting, no additional
pooled-parental MAF\,$\ge 0.15$ filter is applied, since restricting to DI\,$>-25$ already
removes the near-monomorphic loci that filter guards against (it dropped only $\sim$398 of
51{,}612 SNPs). The same statistics were computed at the SNP level for comparison.

\emph{Recombination.} Recombination rate (cM/Mb) was assigned to each marker (and each
unit's representative position) by per-chromosome linear interpolation of the linkage map.
Recombination deciles were defined on the SNP distribution and applied to both levels.

\emph{Panel a.} Within each recombination decile we computed the fraction of differentiated
units (and SNPs) sorted toward each species at $\tau=0.6$ and $\tau=0.8$; grey bars show the
per-decile count of independent eMLG units, deciles with $<100$ units were shaded as
unreliable, and Wilson 95\,\% confidence intervals are shown for eMLG estimates.

\emph{Panel b.} For each $\tau\in\{0.5,0.6,0.7,0.8\}$ and each direction we fitted a logistic
regression of the binary ``sorted toward that species'' indicator on
$\log_{10}(\text{recombination}+0.1)$ across all LD-reduced units. Because units are only
approximately independent, 95\,\% confidence intervals for the recombination coefficient
were obtained from a genomic block bootstrap resampling whole chromosomes (1000 replicates).
The two directions were then compared with a \emph{paired} version of the same bootstrap:
within each replicate the identical resampled chromosomes were used to refit both directions
and their coefficient difference was recorded. The 95\,\% CI of that paired difference and a
two-sided $p$ (twice the fraction of replicates in which the difference reversed sign) test
equality of the two directions' recombination dependence; these $p$-values are printed above
each threshold in panel b.}
{% note
\emph{What the shaded low-recombination deciles actually contain --- counts, not a fraction ---
and why this figure is preliminary.} The three lowest deciles hold 55 units in all; most are
directionally unresolved, and the sorted signal is a handful of large blocks that each collapse
to a single LD-reduced unit. Toward \Fpo{} these are a 350-marker block on chromosome~26, a
190-marker block on chromosome~5 and a 72-marker block on chromosome~6, whose member SNPs
individually reach near-fixation. Toward \Faq{} there is an 89-marker block on chromosome~15
and --- at the more permissive $\tau=0.6$ --- a 384-marker block on chromosome~13, so the
low-recombination \Faq{} signal is \emph{not} confined to tiny clusters. That 384-marker block
is provisional: its pooled-parental MAF ($0.137$) sits just below the $0.15$ floor used by the
genome-wide Module-A sorting, and it enters here only because that gate is deliberately not
applied to the DI25 set (the DI\,$>-25$ selection is itself the gate). Because the
low-recombination contrast rests on so few large units it is sensitive to their handful: it is
significant at $\tau=0.5$--$0.7$ but only marginal at $\tau=0.8$ ($p=0.06$). All of this will be
revisited when the analyses are rerun on the corrected 195-individual sample with reseeded
diagnostic indices.}

%% ==========================================================================
\figblock{fig_size_directional.pdf}
{Figure 3 --- Ancestry sorting direction tracks LD-cluster size, not recombination}
{% caption
\textbf{Among sorted units, direction is governed by LD-cluster size --- and this survives
controlling for recombination.} Unit level ($\tau=0.8$ in a--b; SNPs are all single-marker).
\textbf{(a)} Fraction of units sorted toward \Faq{} (teal) or \Fpo{} (yellow) by cluster-size
class: \Faq{}-directed predominates among small clusters, \Fpo{}-directed among large ones. The
top classes are sparse (32 and 17 units) --- essentially the $\sim$12 large blocks named in
Fig.~2's note (F2233 chr13, F7174 chr26, F6986 chr25, F8626/F8627 chr5, \dots), re-binned by
size. \textbf{(b)} The same relationship without binning: each sorted unit is a point at its
cluster size (log scale), at $y=0$ if sorted toward \Faq{} or $y=1$ toward \Fpo{}, with the
logistic fit $P(\Fpo\mid\text{sorted}) \sim \log_{10}(\text{size})$ and its chromosome
block-bootstrap 95\,\% ribbon. Plotting the units directly shows how sparse the large-cluster
tail is: almost all sorted units are small and \Faq{}-directed, and the rise toward \Fpo{} is
carried by a scatter of a few large blocks --- not an independent axis escaping Fig.~2's
sparsity. \textbf{(c)} A single model, conditional on sorting:
$P(\Fpo\mid\text{sorted}) \sim \log_{10}(\text{recombination}) + \log_{10}(\text{cluster size})$,
fitted across all sorted units per $\tau$ (chromosome block bootstrap, 1000 replicates), on both
the eMLG-consensus sorting (filled) and the representative-marker sorting (open; no averaging
over cluster members). \emph{Cluster size survives and recombination drops out}: the size
coefficient is positive with a 95\,\% CI excluding zero (consensus $+1.1$ at $\tau=0.6$, $+1.6$
at $\tau=0.8$), while the recombination coefficient spans zero at every $\tau$. The size effect
holds on representative markers ($+0.6$ at $\tau=0.6$, $+1.4$ at $\tau=0.8$), so it is not merely
a consensus artefact --- although the consensus does inflate it (roughly halved at $\tau=0.6$).}
{% methods
\emph{Data and units.} The same DI\,$>-25$ units as Figure~2 (no MAF gate), unit level. Cluster
size is the number of diagnostic markers per unit; recombination is the map-interpolated rate
(cM/Mb) at the unit's representative marker. The two are only weakly correlated across units
(Pearson $-0.19$ on the logs).

\emph{Two sorting bases.} \emph{Consensus} scores each unit from its eMLG consensus dosage
(representative SNP for 1--2-marker units), as in Fig.~2; \emph{representative} scores every unit
from its representative marker alone (from the per-SNP analysis), with no averaging over members.

\emph{Panels a--b ($\tau=0.8$).} Panel a bins units by size class (fraction sorted toward each
species, Wilson 95\,\% CIs). Panel b plots each sorted unit at its cluster size with the logistic
fit $P(\Fpo\mid\text{sorted}) \sim \log_{10}(\text{size})$; its ribbon is the same chromosome
block bootstrap as panel c, so b and c share their inference.

\emph{Panel c.} Among units classified as sorted (either direction) we fitted, per
$\tau\in\{0.5,0.6,0.7,0.8\}$ and per basis, a logistic regression of the binary \Fpo{}-vs-\Faq{}
outcome on $\log_{10}(\text{recombination}+0.1)$ and $\log_{10}(\text{cluster size})$ jointly;
95\,\% CIs are from a genomic block bootstrap resampling whole chromosomes (1000 replicates).}
{% note
\emph{Larger clusters sort less often overall --- a measurement effect --- but when they do they
sort \Fpo{}.} The propensity to sort at all declines with cluster size ($P(\text{sorted})$ on
$\log_{10}$ size: $-1.2$, $-1.1$, $-0.8$, $-0.3$ for $\tau=0.5$--$0.8$; CI excludes zero at
$\tau=0.5$--$0.7$): a large eMLG consensus averages over many markers and near-fixes in fewer
populations. Because that averaging could also skew \emph{which} direction survives, panel c
repeats the direction model on representative-marker sorting, which does no averaging; the size
effect persists (smaller), so the size--direction link is not a consensus artefact. This reframes
Figure~2: recombination carried directional signal only because the few large blocks sit in
low-recombination regions --- in the joint model recombination does not survive, and cluster size
does. All preliminary pending the 195-individual, reseeded rerun (see Fig.~2).}

%% ==========================================================================
%% Add the next figure with \figblock{file.pdf}{Title}{caption}{methods}{note-or-empty}

\end{document}
